\uuid{GXM7}
\exo7id{4724}
\titre{exo7 4724}
\auteur{quercia}
\organisation{exo7}
\datecreate{2010-03-16}
\isIndication{false}
\isCorrection{false}
\chapitre{Topologie}
\sousChapitre{Topologie de la droite réelle}
\cpge{mpsi-mp2i.nombres-reels-suites.d, pcsi.nombres-reels-suites.c, ptsi.nombres-reels-suites.c, mpsi-mp2i.fonctions-valeurs-reelles-complexes.d, pcsi.techniques-calcul-fonctions.d, ptsi.techniques-calcul-fonctions.d, bcpst1.suites-reelles, mp-mpi.topologie-espaces-normes.b, pc.espaces-vectoriels-normes.b, psi.espaces-vectoriels-normes.b}
\module{Analyse}
\niveau{L2}
\difficulte{}

\contenu{
\question{
Soit $(u_n)$ une suite r{\'e}elle telle que la suite $(\exp(iu_n))$ converge
et la suite $(|u_{n+1} - u_n|)$ est major{\'e}e par $\alpha < \pi$.
Montrer que $(u_n)$ converge.
}
}
